\documentclass[
	% -- opções da classe memoir --
	12pt,					% tamanho da fonte
	openright,				% capítulos começam em pág ímpar (insere página vazia caso preciso)
	oneside,					% para impressão em verso e anverso. Oposto a oneside
	a4paper,					% tamanho do papel. 
	% -- opções da classe abntex2 --
	%chapter=TITLE,			% títulos de capítulos convertidos em letras maiúsculas
	%section=TITLE,			% títulos de seções convertidos em letras maiúsculas
	%subsection=TITLE,		% títulos de subseções convertidos em letras maiúsculas
	%subsubsection=TITLE,	% títulos de subsubseções convertidos em letras maiúsculas
	% -- opções do pacote babel --
	english,					% idioma adicional para hifenização
	%french,					% idioma adicional para hifenização
	%spanish,				% idioma adicional para hifenização
	brazil					% o último idioma é o principal do documento
	]{abntex2}

% ---------------------
% Pacotes OBRIGATÓRIOS
% ---------------------
\usepackage{lmodern}				% Usa a fonte Latin Modern			
\usepackage[T1]{fontenc}			% Selecao de codigos de fonte.
\usepackage[utf8]{inputenc}		% Codificacao do documento (conversão automática dos acentos)
\usepackage{lastpage}			% Usado pela Ficha catalográfica
\usepackage{indentfirst}			% Indenta o primeiro parágrafo de cada seção.
\usepackage{color}				% Controle das cores
\usepackage{graphicx,graphicx}	% Inclusão de gráficos
\usepackage{epsfig,subfig}		% Inclusão de figuras
\usepackage{microtype} 			% Melhorias de justificação
\usepackage{verbatim}           %comentarios de mais de uma linha
% ---------------------
		
% ---------------------
% Pacotes ADICIONAIS
% ---------------------
\usepackage{lipsum}						% Geração de dummy text
\usepackage{amsmath,amssymb,mathrsfs}	% Comandos matemáticos avançados 
\usepackage{setspace}  					% Para permitir espaçamento simples, 1 1/2 e duplo
\usepackage{tabularx} 					% Para poder ter tabelas com colunas de largura auto-ajustável
\usepackage{afterpage} 					% Para executar um comando depois do fim da página corrente
\usepackage{url} 						% Para formatar URLs (endereços da Web)
\usepackage[table]{xcolor}
\usepackage{float}                      % Para colocar imagens exatamente na posição desejada
\usepackage{mathtools}                  % Símbolos matemáticos
\usepackage{colortbl}                   % Colorir células de tabelas
\usepackage{color}
\usepackage{array}
% ---------------------


% ---------------------
% Pacotes de CITAÇÕES
% ---------------------
\usepackage[brazilian,hyperpageref]{backref}	% Paginas com as citações na bibl
\usepackage[alf]{abntex2cite}				% Citações padrão ABNT (alfa)
%\usepackage[num]{abntex2cite}				% Citações padrão ABNT (numericas)
% ---------------------

% Configurações de CITAÇÕES para abntex2
% --- 
% CONFIGURAÇÕES DE PACOTES
% --- 

% ---
% Configurações do pacote backref
% Usado sem a opção hyperpageref de backref
\renewcommand{\backrefpagesname}{Citado na(s) página(s):~}
% Texto padrão antes do número das páginas
\renewcommand{\backref}{}
% Define os textos da citação
\renewcommand*{\backrefalt}[4]{
	\ifcase #1 %
		Nenhuma citação no texto.%
	\or
		Citado na página #2.%
	\else
		Citado #1 vezes nas páginas #2.%
	\fi}%
% ---

% Inclusão de dados para CAPA e FOLHA DE ROSTO (título, autor, orientador, etc.)
% ---
% Informações de dados para CAPA e FOLHA DE ROSTO
% ---
\titulo{Jogos Sérios para Cibersegurança na Terceira Idade: Uma Revisão Sistemática de Critérios e Recomendações de Design}
\autor{Luis Guilherme Redigolo Crosselli}
\local{Santo André - SP}
\data{Abril de 2025}
\orientador{Denise Hideko Goya}
%\coorientador{Fulano Nome do Coorientador}
\instituicao{%
  Universidade Federal do ABC -- UFABC
  \par
  Centro de Matemática, Computação \& Cognição 
  \par
  Bacharelado em Ciência da Computação}
\tipotrabalho{Projeto de Graduação em Computação}
% O preambulo deve conter o tipo do trabalho, o objetivo,
% o nome da instituição e a área de concentração
\preambulo{\textbf{Projeto de Graduação em Computação} apresentado como parte dos requisitos necessários para a obtenção do Título de Bacharel em Ciência da Computação.}
% ---

% Inclui Configurações de aparência do PDF Final
%  Configurações de aparência do PDF final
% NÃO ALTERAR!!!

% alterando o aspecto da cor azul
\definecolor{blue}{RGB}{41,5,195}

% informações do PDF
\makeatletter
\hypersetup{
     	%pagebackref=true,
		pdftitle={\@title}, 
		pdfauthor={\@author},
    		pdfsubject={\imprimirpreambulo},
	    pdfcreator={LaTeX with abnTeX2},
		pdfkeywords={abnt}{latex}{abntex}{abntex2}{trabalho acadêmico}, 
		colorlinks=true,       		% false: boxed links; true: colored links
    		linkcolor=blue,          	% color of internal links
    		citecolor=blue,        		% color of links to bibliography
    		filecolor=magenta,      		% color of file links
		urlcolor=blue,
		bookmarksdepth=4
} 
\makeatother
% --- 

% O tamanho da identação do parágrafo é dado por:
\setlength{\parindent}{1.3cm}

% Controle do espaçamento entre um parágrafo e outro:
\setlength{\parskip}{0.2cm}  % tente também \onelineskip

% ---------------------
% Compila o indice
% ---------------------
\makeindex
% ---------------------

%%%%%%%%%%%%%%%%%%%%%%%%%%%
%%  INICIO DO DOCUMENTO  %%
%%%%%%%%%%%%%%%%%%%%%%%%%%%
\begin{document}

% Retira espaço extra obsoleto entre as frases.
\frenchspacing

% ----------------------------------------------------------
% ELEMENTOS PRÉ-TEXTUAIS (Capa, Resumo, Abstract, etc.)
% ----------------------------------------------------------
\pretextual

% Capa
% ---
% Impressão da Capa
% ---
  \begin{capa}%
    \begin{figure}[h!]%
        \centering%
        \includegraphics[scale=1.2]{figs/logo.png}%
      \end{figure}%
    \center
	\ABNTEXchapterfont\large{Universidade Federal do ABC \\ Centro de Matemática, Computação \& Cognição \\ Bacharelado em Ciência da Computação}
	%\vspace{1.5cm}

    \vfill
    \ABNTEXchapterfont\bfseries\LARGE\imprimirtitulo
    \vfill

	%\vfill
	\ABNTEXchapterfont\large\imprimirautor
	\vfill
%
	
	
    \large\imprimirlocal, \large\imprimirdata

    \vspace*{1cm}
  \end{capa}
% ---

% Folha de rosto (o * indica que haverá a ficha bibliográfica)
\imprimirfolhaderosto*

 \begin{comment}
% Imprimir Ficha Catalografica
% ---
% Ficha Catalográfica
% ---
% Isto é um exemplo de Ficha Catalográfica, ou ``Dados internacionais de
% catalogação-na-publicação''. Você pode utilizar este modelo como referência. 
% Porém, talvez a biblioteca lhe fornece um PDF
% com a ficha catalográfica definitiva após a defesa do trabalho. Quando estiver
% com o documento, salve-o como PDF no diretório do seu projeto e substitua todo
% o conteúdo de implementação deste arquivo pelo comando abaixo:
%
% \begin{fichacatalografica}
%     \includepdf{fig_ficha_catalografica.pdf}
% \end{fichacatalografica}
\begin{fichacatalografica}
	\vspace*{\fill}					% Posição vertical
	\hrule							% Linha horizontal
	\begin{center}					% Minipage Centralizado
	\begin{minipage}[c]{12.5cm}		% Largura
	
	\imprimirautor
	
	\hspace{0.5cm} \imprimirtitulo  / \imprimirautor. --
	\imprimirlocal, \imprimirdata-
	
	\hspace{0.5cm} \pageref{LastPage} p. : il. (algumas color.) ; 30 cm.\\
	
	\hspace{0.5cm} \imprimirorientadorRotulo~\imprimirorientador\\
	
	\hspace{0.5cm}
	\parbox[t]{\textwidth}{\imprimirtipotrabalho~--~\imprimirinstituicao,
	\imprimirdata.}\\
	
	\hspace{0.5cm}
		1. Palavra-chave1.
		2. Palavra-chave2.
		I. Orientador.
		II. Universidade xxx.
		III. Faculdade de xxx.
		IV. Título\\ 			
	
	\hspace{8.75cm} CDU 02:141:005.7\\
	
	\end{minipage}
	\end{center}
	\hrule
\end{fichacatalografica}
% ---

% Inserir Folha de Aprovação
% ---
% Assinaturas
% ---
% Isto é um exemplo de Folha de aprovação, elemento obrigatório da NBR
% 14724/2011 (seção 4.2.1.3). Você pode utilizar este modelo até a aprovação
% do trabalho. Após isso, substitua todo o conteúdo deste arquivo por uma
% imagem da página assinada pela banca com o comando abaixo:
%
% \includepdf{folhadeaprovacao_final.pdf}
%
\begin{folhadeaprovacao}

  \begin{center}
    {\ABNTEXchapterfont\large\imprimirautor}

    \vspace*{\fill}\vspace*{\fill}
    \begin{center}
      \ABNTEXchapterfont\bfseries\Large\imprimirtitulo
    \end{center}
    \vspace*{\fill}
    
    \hspace{.45\textwidth}
    \begin{minipage}{.5\textwidth}
        \imprimirpreambulo
    \end{minipage}%
    \vspace*{\fill}
   \end{center}
        
 % Isso na versao final do trabalho!!!       
   Trabalho aprovado. \imprimirlocal, 01 de janeiro de 2014:

   \assinatura{\textbf{\imprimirorientador} \\ Orientador} 

   \assinatura{\textbf{Professor} \\ Convidado 1}
   \assinatura{\textbf{Professor} \\ Convidado 2}
   
      
   \begin{center}
    \vspace*{0.5cm}
    {\large\imprimirlocal}
    \par
    {\large\imprimirdata}
    \vspace*{1cm}
  \end{center}
  
\end{folhadeaprovacao}
% ---

% Dedicatória
% ---
% Dedicatória
% ---
\begin{dedicatoria}
   \vspace*{\fill}
   \centering
   \noindent
   \textit{ Aos verme que roeu as frias carnes de meu cadáver.} \vspace*{\fill}
\end{dedicatoria}
% ---

% Agradecimentos
% ---
% Agradecimentos
% ---
\begin{agradecimentos}



Agradeço a Xuxa, meus pais, cachorro, gato e papagaio, por ...

Agradeço ao meu orientador, XXXXXXXXX, por todos os conselhos, pela paciência e ajuda nesse período.

Aos meus amigos ...

Aos professores ...

À XXXXXX pelo apoio financeiro para realização deste trabalho de pesquisa.

\end{agradecimentos}
%% ---

% Epígrafe
% ---
% Epígrafe
% ---
\begin{epigrafe}
    \vspace*{\fill}
	\begin{flushright}
		\textit{``Não sei o que, \\
		          não sei o que,\\
                  não sei o que lá.''\\
		          (Autor Desconhecido)}
	\end{flushright}
\end{epigrafe}
% ---

% Resumo e Abstract
% ---
% RESUMOS
% ---

% RESUMO em português
\setlength{\absparsep}{18pt} % ajusta o espaçamento dos parágrafos do resumo
\begin{resumo}

O presente trabalho realiza uma Revisão Sistemática da Literatura (RSL) com o objetivo de identificar critérios, recomendações e metodologias aplicadas no desenvolvimento de jogos sérios voltados à população idosa, com foco especial em cibersegurança. O envelhecimento populacional e a crescente inserção digital dos idosos tornam esse público mais vulnerável a ameaças online, reforçando a importância de soluções educacionais acessíveis e eficazes. A pesquisa analisou 62 estudos relevantes, destacando lacunas no uso de diretrizes específicas de segurança digital nos jogos existentes. Além disso, observou-se uma carência de abordagens participativas no levantamento de requisitos. Como contribuição, o estudo propõe a consolidação de recomendações que sirvam como guia prático para desenvolvedores e pesquisadores interessados em promover a inclusão digital segura da população idosa por meio de jogos sérios.


 \textbf{Palavras-chaves}: Jogos Sérios, Cibersegurança, Pessoas Idosas, Design Centrado no Usuário, Inclusão Digital.
\end{resumo}

% ABSTRACT in english
\begin{resumo}[Abstract]
 \begin{otherlanguage*}{english}
    This work presents a Systematic Literature Review (SLR) aiming to identify criteria, requirements, and methodologies applied in the development of serious games for the elderly population, with a special focus on cybersecurity. The aging of the population and the increasing digital inclusion of older adults make this group more vulnerable to online threats, highlighting the need for accessible and effective educational solutions. The review analyzed 62 relevant studies, revealing gaps in the application of specific digital security guidelines in existing games. Furthermore, a lack of participatory approaches in requirements elicitation was identified. As a contribution, the study proposes a set of consolidated recommendations to serve as a practical guide for developers and researchers interested in promoting safe digital inclusion for the elderly through serious games.

   \vspace{\onelineskip}
 
   \noindent 
   \textbf{Keywords}: Serious Games, Cybersecurity, Older Adults, User-Centered Design, Digital Inclusion.
 \end{otherlanguage*}
\end{resumo}

% Lista de ilustrações
\pdfbookmark[0]{\listfigurename}{lof}
\listoffigures*
\cleardoublepage

% Lista de tabelas
\pdfbookmark[0]{\listtablename}{lot}
\listoftables
\cleardoublepage

% Lista de abreviaturas e siglas
\begin{siglas}
  \item[ABNT] Associação Brasileira de %Normas Técnicas
  \item[abnTeX] Normas para TeX
\end{siglas}

% Lista de símbolos
\begin{simbolos}
  \item[$ \Gamma $] Letra grega Gama
  \item[$ \Lambda $] Lambda
  \item[$ \zeta $] Letra grega minúscula zeta
  \item[$ \in $] Pertence
\end{simbolos}

% Inserir o SUMÁRIO
\pdfbookmark[0]{\contentsname}{toc}
\tableofcontents
\cleardoublepage
\end{comment} 

% Inserir Folha de Aprovação
% % ---
% Assinaturas
% ---
% Isto é um exemplo de Folha de aprovação, elemento obrigatório da NBR
% 14724/2011 (seção 4.2.1.3). Você pode utilizar este modelo até a aprovação
% do trabalho. Após isso, substitua todo o conteúdo deste arquivo por uma
% imagem da página assinada pela banca com o comando abaixo:
%
% \includepdf{folhadeaprovacao_final.pdf}
%
\begin{folhadeaprovacao}

  \begin{center}
    {\ABNTEXchapterfont\large\imprimirautor}

    \vspace*{\fill}\vspace*{\fill}
    \begin{center}
      \ABNTEXchapterfont\bfseries\Large\imprimirtitulo
    \end{center}
    \vspace*{\fill}
    
    \hspace{.45\textwidth}
    \begin{minipage}{.5\textwidth}
        \imprimirpreambulo
    \end{minipage}%
    \vspace*{\fill}
   \end{center}
        
 % Isso na versao final do trabalho!!!       
   Trabalho aprovado. \imprimirlocal, 01 de janeiro de 2014:

   \assinatura{\textbf{\imprimirorientador} \\ Orientador} 

   \assinatura{\textbf{Professor} \\ Convidado 1}
   \assinatura{\textbf{Professor} \\ Convidado 2}
   
      
   \begin{center}
    \vspace*{0.5cm}
    {\large\imprimirlocal}
    \par
    {\large\imprimirdata}
    \vspace*{1cm}
  \end{center}
  
\end{folhadeaprovacao}
% ---

% Resumo e Abstract
% ---
% RESUMOS
% ---

% RESUMO em português
\setlength{\absparsep}{18pt} % ajusta o espaçamento dos parágrafos do resumo
\begin{resumo}

O presente trabalho realiza uma Revisão Sistemática da Literatura (RSL) com o objetivo de identificar critérios, recomendações e metodologias aplicadas no desenvolvimento de jogos sérios voltados à população idosa, com foco especial em cibersegurança. O envelhecimento populacional e a crescente inserção digital dos idosos tornam esse público mais vulnerável a ameaças online, reforçando a importância de soluções educacionais acessíveis e eficazes. A pesquisa analisou 62 estudos relevantes, destacando lacunas no uso de diretrizes específicas de segurança digital nos jogos existentes. Além disso, observou-se uma carência de abordagens participativas no levantamento de requisitos. Como contribuição, o estudo propõe a consolidação de recomendações que sirvam como guia prático para desenvolvedores e pesquisadores interessados em promover a inclusão digital segura da população idosa por meio de jogos sérios.


 \textbf{Palavras-chaves}: Jogos Sérios, Cibersegurança, Pessoas Idosas, Design Centrado no Usuário, Inclusão Digital.
\end{resumo}

% ABSTRACT in english
\begin{resumo}[Abstract]
 \begin{otherlanguage*}{english}
    This work presents a Systematic Literature Review (SLR) aiming to identify criteria, requirements, and methodologies applied in the development of serious games for the elderly population, with a special focus on cybersecurity. The aging of the population and the increasing digital inclusion of older adults make this group more vulnerable to online threats, highlighting the need for accessible and effective educational solutions. The review analyzed 62 relevant studies, revealing gaps in the application of specific digital security guidelines in existing games. Furthermore, a lack of participatory approaches in requirements elicitation was identified. As a contribution, the study proposes a set of consolidated recommendations to serve as a practical guide for developers and researchers interested in promoting safe digital inclusion for the elderly through serious games.

   \vspace{\onelineskip}
 
   \noindent 
   \textbf{Keywords}: Serious Games, Cybersecurity, Older Adults, User-Centered Design, Digital Inclusion.
 \end{otherlanguage*}
\end{resumo}

% Inserir o SUMÁRIO
\pdfbookmark[0]{\contentsname}{toc}
\tableofcontents*
\cleardoublepage



% ----------------------------------------------------------
% ELEMENTOS TEXTUAIS (Capítulos)
% ----------------------------------------------------------
\textual
% Elementos textuais com numeração arábica
\pagenumbering{arabic}
% Reinicia a contagem do número de páginas
\setcounter{page}{1}


% Inclui cada capitulo do Trabalho
% ----------------------------------------------------------
% Introdução 
% Capítulo sem numeração, mas presente no Sumário
% ----------------------------------------------------------

\chapter[Introdução]{Introdução}
\label{cap:introducao}
% \addcontentsline{toc}{chapter}{Introdução}



 

% Demonstração de citação: o software de análise foi desenvolvido em Python~\cite{van1995python}, utilizando as bibliotecas Pandas~\cite{mckinney2010data} e Scikit-learn~\cite{scikit-learn}.

O envelhecimento da população é um fenômeno global que traz consigo novas demandas sociais e tecnológicas. Com a crescente digitalização de serviços e interações, a população idosa tem se tornado cada vez mais conectada, mas também mais vulnerável a ameaças cibernéticas~\cite{morrison2021how}. Nesse contexto, os jogos sérios (serious games) emergem como uma ferramenta promissora para educação e treinamento, oferecendo um ambiente interativo e motivador para o desenvolvimento de habilidades ~\cite{connolly2012systematic}. Especificamente na área de cibersegurança, os jogos sérios podem desempenhar um papel crucial na capacitação de idosos para navegar no mundo digital com mais segurança.

Embora os benefícios dos jogos sérios para a população idosa sejam reconhecidos, abrangendo melhorias cognitivas, físicas e sociais ~\cite{konstantinidis2010integration, mccallum2012gamification}, existe uma lacuna significativa no que diz respeito a jogos focados em cibersegurança e projetados especificamente para este público. Muitos jogos existentes não consideram adequadamente as necessidades, preferências e limitações dos usuários idosos, comprometendo sua eficácia e engajamento ~\cite{mol2021desirable}. A aplicação de princípios de Desenho Universal e a atenção à acessibilidade são fundamentais, mas ainda pouco exploradas nesse domínio específico ~\cite{belarmino2021criterios}.

O desenvolvimento de jogos sérios eficazes para idosos apresenta desafios únicos, especialmente na fase de levantamento de requisitos ~\cite{beristain2021standardizing}. Garantir que o jogo seja não apenas funcional, mas também acessível, intuitivo e motivador para o público-alvo exige uma compreensão profunda de suas características e expectativas. Metodologias de design centradas no usuário (\textit{User-Centered Design} - UCD) são essenciais ~\cite{brox2017user}, porém, a adaptação de técnicas de elicitação de requisitos para considerar as particularidades cognitivas, físicas e a familiaridade com tecnologia dos idosos ainda é um campo que necessita de maior investigação ~\cite{fua2013designing}.

Diante desse cenário, identifica-se a ausência de um conjunto consolidado de diretrizes e critérios, fundamentado em evidências, que possa guiar o desenvolvimento de jogos sérios sobre cibersegurança para a população idosa. Questões fundamentais permanecem em aberto: Quais são as diretrizes de design essenciais para esses jogos? Como eles devem ser validados junto ao público-alvo? E quais abordagens de levantamento de requisitos se mostram mais eficazes para capturar as necessidades desse grupo específico?

Este trabalho tem como objetivo principal preencher essa lacuna por meio de uma Revisão Sistemática da Literatura (RSL). Procura-se identificar, analisar e sintetizar os critérios, diretrizes, boas práticas e desafios relacionados ao desenvolvimento e avaliação de jogos sérios voltados para pessoas idosas, especialmente no contexto de cibersegurança. Como resultado, pretende-se consolidar um conjunto de recomendações que possa servir como um guia prático para desenvolvedores, designers e pesquisadores interessados em criar soluções de jogos mais eficazes e acessíveis para este público.

A metodologia adotada seguirá um protocolo de RSL, envolvendo buscas sistemáticas em bases de dados acadêmicas relevantes (como \textit{IEEE Xplore}, \textit{ACM Digital Library}, \textit{Web of Science} e \textit{SBC Open Lib}), aplicação de critérios de inclusão e exclusão pré-definidos, e análise qualitativa dos estudos selecionados. Espera-se que a contribuição deste trabalho resida na consolidação do conhecimento existente e na oferta de um referencial prático e baseado em evidências para o avanço da área de jogos sérios para a cibersegurança e inclusão digital da população idosa.

\section{Justificativa}
%\section{Motivação}
A vulnerabilidade crescente de pessoas idosas a ameaças cibernéticas, combinada com o potencial dos jogos sérios como ferramenta educacional engajadora, justifica a necessidade de investigar e consolidar recomendações específicas para o desenvolvimento de tais jogos. A falta de diretrizes claras dificulta a criação de soluções eficazes e acessíveis, reforçando a importância desta pesquisa para orientar futuros desenvolvimentos na área.

\section{Objetivos}

\subsection{Objetivo Geral}
Realizar uma Revisão Sistemática da Literatura (RSL) para identificar e analisar os critérios, diretrizes e metodologias de levantamento de requisitos propostos no desenvolvimento de jogos sérios voltados para o público idoso, com foco em aplicações relacionadas à cibersegurança e segurança digital.

\subsection{Objetivos Específicos}
Os objetivos específicos são:
\begin{itemize}
    \item Apresentar o protocolo da Revisão Sistemática da Literatura (RSL) realizada, incluindo as questões de pesquisa, \textit{strings} de busca, bases de dados consultadas, e critérios de inclusão e exclusão.
    \item Descrever o processo de seleção dos artigos relevantes.
    \item Apresentar os resultados preliminares da RSL, com foco na identificação dos artigos selecionados que abordam recomendações para jogos sérios para idosos, destacando aqueles relacionados à cibersegurança.
    \item Discutir a metodologia de trabalho e o cronograma previsto para as próximas fases do projeto (PGC-II e PGC-III), que incluirão a análise aprofundada dos artigos selecionados, e a consolidação dos resultados.
\end{itemize}

%sugestão
\section{Estrutura da Monografia}
O presente relatório está estruturado da seguinte forma: o Capítulo~\ref{cap:fundamentacao_teorica} apresenta a fundamentação teórica sobre jogos sérios, o público idoso e tecnologia, e cibersegurança. O Capítulo~\ref{cap:trabalhos_relacionados} discute trabalhos relacionados, incluindo outras revisões de literatura. O Capítulo~\ref{cap:metodologia} detalha a metodologia da RSL e o cronograma do projeto. O Capítulo~\ref{cap:resultados} apresenta os resultados preliminares da busca e seleção de artigos. Finalmente, o Capítulo~\ref{cap:conclusoes} sumariza os achados desta fase inicial e delineia os próximos passos da pesquisa.

% SUGESTAO: Exemplo de estrutura da monografia:

% O presente relatório está estruturado da seguinte forma: o Capítulo~\ref{cap:fundamentacao_teorica} apresenta..., o Capítulo~\ref{cap:trabalhos_relacionados}... O Capítulo~\ref{cap:metodologia} ..., o Capítulo~\ref{cap:resultados} .... Finalmente, o Capítulo~\ref{cap:conclusao} apresenta ...




%\include{capitulos/justificativa}
%\include{capitulos/objetivos}
\chapter{Fundamentação Teórica}
\label{cap:fundamentacao_teorica}

A seguir, será apresentada a definição e o escopo dos jogos sérios, além de suas distinções com sistemas gamificados, essenciais para entender o impacto desses jogos no contexto da saúde e bem-estar dos idosos. Outros conceitos como o Design Centrado no Usuário também são introduzidos.

\section{Jogos Sérios}

Jogos sérios são aplicações digitais que combinam elementos de entretenimento e multimídia com o propósito explícito de promover experiências voltadas para objetivos além do divertimento. Esses jogos engajam os usuários por meio de mecânicas lúdicas, mas diferenciam-se dos jogos puramente recreativos por incorporarem conteúdos educacionais, instrucionais, terapêuticos ou de treinamento. A definição mais amplamente aceita caracteriza jogos sérios como aqueles cujo propósito principal não é o entretenimento, ainda que este seja utilizado como meio para facilitar a assimilação de conhecimentos, habilidades ou valores. Assim, o termo “sério” refere-se à intenção de transmitir ao jogador algum tipo de conteúdo significativo, oriundo de experiências ou saberes aplicáveis a contextos reais como saúde, educação ou comunicação interpessoal~\cite{laamarti2014overview}.

A literatura destaca que os jogos sérios são caracterizados por três componentes fundamentais: experiência, entretenimento e multimídia~\cite{laamarti2014overview}. A experiência diz respeito à vivência do jogador com o conteúdo proposto, o entretenimento atua como fator de engajamento e motivação, e os recursos multimídia abrangem os canais sensoriais utilizados na interação (visuais, auditivos, táteis etc.). Diferentes perspectivas acadêmicas e industriais concordam que o sucesso desses jogos reside na capacidade de equilibrar o conteúdo sério com uma experiência envolvente e interativa~\cite{nazry2017mood}. Dessa forma, jogos sérios constituem uma ferramenta promissora para impactar positivamente diversas áreas da sociedade, desde a educação formal até o treinamento militar e a reabilitação médica~\cite{michael2005serious, canapa2025interactive}.

\section{Gamificação vs Jogos Sérios}

É importante distinguir jogos sérios de sistemas gamificados. A \textit{gamificação} refere-se à aplicação de elementos de jogos (como pontos, rankings e recompensas) em contextos não lúdicos, como aplicativos de saúde ou educação \cite{deterding2011gamefulness}. Já os \textit{jogos sérios} são sistemas completos, com mecânicas próprias e objetivos definidos. Para idosos, os jogos sérios oferecem experiências mais imersivas e significativas do que a simples gamificação de tarefas.

\section{Design Centrado no Usuário Idoso}

O desenvolvimento de jogos para idosos exige atenção especial às diretrizes de design centrado no usuário \cite{abras2004user}, priorizando \textit{usabilidade} e \textit{acessibilidade}. Barreiras comuns enfrentadas por idosos incluem visão reduzida, dificuldades auditivas, lentidão na resposta motora e declínio cognitivo. Portanto, de acordo com o W3C \cite{w3c2024wcag}, interfaces devem ser simples, apresentar textos legíveis, ter feedbacks claros e tempos de resposta ajustáveis.

Segundo \cite{fisk2009designing}, é essencial considerar fatores como familiaridade com tecnologia, ergonomia do dispositivo e linguagem acessível. Outro ponto importante é oferecer autonomia e controle ao usuário, respeitando seu ritmo e suas limitações.

\section{Aspectos do Envelhecimento e Interação Digital}

Com o envelhecimento, é comum a ocorrência de declínio cognitivo (como redução da memória de curto prazo e atenção), bem como diminuição da acuidade visual e coordenação motora \cite{charness2009aging}. Estudos mostram que jogos digitais podem atenuar esses efeitos ao estimular áreas específicas do cérebro e manter os usuários engajados cognitivamente \cite{green2008plasticity}.

Além disso, jogos podem servir como ferramentas de socialização, ajudando a combater o isolamento social — um fator de risco para doenças mentais e declínio funcional em idosos \cite{cotten2012internet}.

\section{Tecnologias Aplicadas em Jogos para Idosos}

Diversas tecnologias têm sido utilizadas no desenvolvimento de jogos sérios voltados para o público idoso. Tablets e smartphones com interfaces sensíveis ao toque são preferidos por sua simplicidade de uso \cite{sayago2011everyday}. Tecnologias como reconhecimento de voz, sensores de movimento e realidade aumentada também são promissoras, pois oferecem formas naturais de interação \cite{gerling2012game}. No entanto, o uso dessas tecnologias exige cuidados com acessibilidade e custo, além do esforço necessário para ajustar os equipamentos ou preparar o ambiente em que o jogo será executado.

\chapter{Trabalhos Relacionados}\label{cap:trabalhos_relacionados}

A relação entre jogos digitais e segurança da informação no contexto da terceira idade tem sido explorada por diferentes estudos, com enfoques variados em alfabetização digital, reabilitação cognitiva, usabilidade e proteção contra golpes virtuais. Esta seção apresenta os principais trabalhos identificados que abordam tais temas, com destaque para suas propostas, metodologias e contribuições.

Um estudo relevante nesse contexto é o de Pillon e Silva \cite{pillon2022checklist}, que propuseram um checklist estruturado com 30 itens organizados nas categorias de terapia, motivação, interação e segurança, com o objetivo de orientar o desenvolvimento e a avaliação de jogos digitais em realidade virtual para a reabilitação de idosos. O checklist foi fundamentado em uma revisão sistemática de literatura e validado por meio de entrevistas semiestruturadas com especialistas, sendo uma importante ferramenta metodológica para garantir a adequação dos jogos às necessidades do público sênior.

O jogo Quiz60+ foi desenvolvido com o objetivo de instruir idosos sobre os principais golpes virtuais, incluindo fraudes via WhatsApp, sites falsos e mensagens armadilhas. Seu conteúdo foi estruturado em formato de quiz com múltipla escolha, oferecendo feedback imediato e interface intuitiva. A validação envolveu um pequeno grupo de participantes que testaram o jogo em dispositivos Android e responderam a um questionário de avaliação \cite{amorim2024quiz}.

Outro trabalho relevante é o de \cite{scarpioni2016desenvolvimento}, que propôs ambientes simulados de e-mail e redes sociais com foco na identificação de riscos durante a navegação. Apesar de não constituírem jogos no sentido tradicional, os simuladores visam promover o aprendizado por meio de interação e reflexão sobre comportamentos digitais inseguros.

Além dos jogos com foco em segurança digital, outros trabalhos investigam o uso de jogos sérios na reabilitação cognitiva e física de idosos. Tais estudos, embora não abordem diretamente a cibersegurança, oferecem contribuições importantes sobre usabilidade, motivação e adaptação de interfaces para esse público.

Um levantamento sistemático de grande abrangência foi realizado por \cite{canapa2025interactive}, que analisaram 113 estudos sobre jogos sérios interativos voltados ao treinamento cognitivo de adultos mais velhos. Os autores examinaram dimensões como as tecnologias interativas utilizadas (com destaque para tablets, VR/AR e computadores), os tipos de atividades propostas (como quebra-cabeças, simuladores de tarefas cotidianas e jogos de memória), os perfis dos usuários (incluindo idosos saudáveis, com comprometimento cognitivo leve ou demência) e os efeitos dessas intervenções sobre funções cognitivas específicas como memória, atenção e funções executivas. O trabalho destaca a predominância do uso de tablets, a recorrência de atividades baseadas em correspondência de pares e simulações do cotidiano, e o uso ainda incipiente de tecnologias imersivas. Apesar de seu escopo abrangente, o estudo se concentra majoritariamente em aspectos técnicos e clínicos dos jogos sérios para reabilitação cognitiva, não abordando diretamente os desafios relacionados à segurança digital ou à alfabetização em cibersegurança do público idoso.

Dentre os trabalhos relacionados, destacam-se revisões de literatura que apresentam objetivos semelhantes aos deste estudo, embora divirjam em escopo ou foco.

\cite{tong2014cognitive} explora o desenvolvimento de um jogo sério voltado à avaliação cognitiva de idosos, com foco particular na detecção de condições transitórias como o delirium — um desafio clínico de alta relevância, mas frequentemente subdiagnosticado. O trabalho descreve o processo iterativo de prototipação envolvendo diferentes plataformas (\textit{Wiimote}, \textit{Kinect} e \textit{tablets}) até a escolha final por um jogo de acertar alvos em dispositivos com tela sensível ao toque. A escolha baseou-se em sua portabilidade, baixo custo e maior acessibilidade para pacientes idosos. O jogo foi projetado para medir a capacidade de inibição, uma função executiva frequentemente comprometida com o envelhecimento, e passou por testes de usabilidade com adultos e idosos. Os resultados evidenciaram correlações significativas entre o desempenho no jogo e testes cognitivos padronizados, embora tenham sido apontadas limitações ergonômicas da interface para o público idoso.

\cite{ning2020review} apresentam uma revisão abrangente sobre o uso de jogos sérios no cuidado de pessoas com demência, classificando os jogos conforme os sintomas em diferentes estágios da doença e propondo um modelo de avaliação baseado em múltiplos métodos e grupos. Embora o foco em jogos voltados à população idosa seja um ponto de convergência com este trabalho, a abordagem dos autores centra-se na reabilitação cognitiva e física de doenças neurodegenerativas, sem diretamente elencar recomendações. Diferentemente disso, esta pesquisa sistematiza critérios e recomendações de design voltados à criação de jogos sérios que auxiliem a inclusão digital segura de idosos, com ênfase na proteção contra ameaças digitais.

\cite{lee2019study} exploram os fatores de design de jogos sérios voltados à população idosa, com ênfase na interseção entre propósito funcional e elementos de diversão. O trabalho propõe um modelo conceitual para guiar o design de jogos para esse público, destacando a necessidade de superar barreiras culturais e técnicas na Coreia do Sul. Apesar da relevância ao discutir características específicas do público sênior, o artigo não trata de domínios como a cibersegurança, tampouco aborda metodologias de levantamento de requisitos centradas no usuário idoso.

\cite{cedillo2025systematic} realizam uma revisão sistemática sobre tecnologias de interação humano-computador aplicadas à saúde física e cognitiva de idosos, identificando dispositivos de captura de movimento e abordagens multimodais de interação. Embora também utilizem uma abordagem sistemática e abordem o uso de jogos sérios e exergames, seu foco está na tecnologia assistiva para reabilitação e suporte à saúde, não considerando aplicações educacionais no contexto de segurança digital. Em contraste, o presente estudo propõe diretrizes específicas para o design de jogos educativos sobre cibersegurança voltados a esse público, considerando limitações cognitivas e sensoriais decorrentes do envelhecimento.

\cite{postal2019user} oferecem uma revisão sistemática sobre métodos de avaliação de interfaces voltadas a idosos, com foco em aplicações de realidade virtual. Os autores destacam a falta de metodologias consolidadas para avaliação de usabilidade em interfaces tridimensionais destinadas a esse público. Embora compartilhe com este trabalho a preocupação com acessibilidade e adequação de tecnologias interativas ao perfil da população idosa, a proposta de Postal e Rieder está centrada na avaliação de usabilidade, enquanto este estudo foca na definição de critérios de design e conteúdo para jogos sérios.

\cite{bassano2022visualization} revisam tecnologias de visualização e interação utilizadas em jogos sérios e exergames para avaliação e treinamento cognitivo. O estudo destaca o uso crescente de realidade virtual e aumentada e discute a validação das soluções propostas, mas observa limitações na adoção clínica e na disponibilidade de dados.

De modo geral, os trabalhos analisados revelam uma preocupação crescente com o treinamento cognitivo, a usabilidade e o letramento digital de idosos. No entanto, observa-se uma lacuna significativa na sistematização de critérios e recomendações que integrem, de forma estruturada, princípios de segurança da informação ao design de jogos sérios voltados a esse público. Embora alguns estudos abordem indiretamente temas relacionados à proteção digital, poucos exploram a cibersegurança como eixo central no desenvolvimento das soluções. Além disso, é raro encontrar abordagens que combinem metodologias de levantamento participativo de requisitos com diretrizes de acessibilidade e princípios do desenho universal, de modo a contemplar de forma inclusiva as diferentes necessidades cognitivas, motoras e tecnológicas do público idoso.

Este trabalho diferencia-se dos citados anteriormente nos seguintes aspectos: (i) realiza uma análise específica e sistemática da interseção entre segurança digital e o design de jogos sérios voltados ao público idoso — uma abordagem ainda pouco explorada e consolidada na literatura; (ii) propõe a consolidação de um conjunto de critérios e diretrizes organizados em forma de checklist, com o objetivo de orientar o desenvolvimento de jogos sérios que sejam acessíveis, engajadores e alinhados às práticas de proteção digital; e (iii) adota uma abordagem que vai além do uso de jogos sérios como ferramentas exclusivamente voltadas ao treinamento cognitivo ou à prevenção de doenças neurodegenerativas.

\chapter{Metodologia}
\label{cap:metodologia}

Para identificar os critérios e recomendações fundamentais para o desenvolvimento de jogos sérios sobre cibersegurança para pessoas idosas, foi adotada a metodologia de Revisão Sistemática da Literatura (RSL). Esta abordagem permite uma investigação científica rigorosa e reprodutível da literatura existente, seguindo um protocolo bem definido, conforme descrito por \cite{kitchenham2004procedures}.

\section{Protocolo da RSL}\label{sec:protocolo_rsl}

\subsection{Questões de Pesquisa}\label{subsec:questoes_pesquisa}

As seguintes questões de pesquisa foram definidas para nortear a RSL:
\begin{enumerate}
    \item Quais as recomendações, requisitos, e critérios identificados no desenvolvimento de jogos para o público idoso?
    \item Como jogos desenvolvidos tendo o público idoso como alvo são validados e testados?
    \item Que abordagens são usadas para levantamento de requisitos nestes jogos e quais produzem melhores resultados?
\end{enumerate}

\subsection{Critérios de Inclusão e Exclusão}\label{subsec:criterios}

Para garantir a relevância e qualidade dos estudos selecionados, foram estabelecidos os seguintes critérios:

\textbf{Critérios de inclusão:}
\begin{enumerate}
    \item Apresenta recomendações específicas voltados ao público idoso em contextos de jogos ou tecnologias interativas;
    \item Discute metodologias de levantamento de requisitos com participação do público idoso ou voltadas a suas características;
    \item Inclui avaliação empírica ou reflexões sobre a efetividade dos critérios adotados para o público idoso;
    \item Relaciona as recomendações e critérios com aspectos de cibersegurança ou proteção digital.
\end{enumerate}

\textbf{Critérios de exclusão:}
\begin{enumerate}
    \item Trabalhos que não sejam artigos de pesquisa;
    \item Trabalhos que não estejam escritos em Língua Portuguesa ou Inglesa;
    \item Trabalhos duplicados ou que descrevam o mesmo projeto;
    \item Trabalhos publicados antes de 2015;
    \item Trabalhos que mencionam o público idoso apenas de forma tangencial, sem foco na proposta de design ou avaliação de jogos para esse público;
    \item Trabalhos que fazem uso de sensores de movimento como principal característica;
    \item Revisões de literatura;
    \item Trabalhos não disponíveis para acesso;
\end{enumerate}

\subsection{Bases de Dados e Estratégia de Busca}\label{subsec:bases_busca}

A busca foi realizada nas seguintes bases de dados:
\begin{itemize}
    \item IEEE Xplore
    \item ACM Digital Library
    \item Web of Science
    \item SBC OpenLib (para publicações nacionais)
\end{itemize}

As bases de dados foram escolhidas por sua ampla relevância e abrangência na área de Ciência da Computação, Engenharia e Tecnologia da Informação. A escolha dessas bases visa garantir a abrangência, diversidade e confiabilidade dos estudos coletados e são amplamente utilizadas em revisões sistemáticas e pesquisas científicas. Para cada base, foram utilizadas strings de busca adaptadas à sua sintaxe específica, mas mantendo a mesma lógica de combinação de termos relacionados a:

\begin{enumerate}
    \item Jogos: \textit{game, games, gaming, jogo, jogos}
    \item Educacional/Sério: \textit{educational, serious, educacional, sério}
    \item Público-alvo: \textit{elderly, old adult, old adults, senior, idoso, terceira idade}
\end{enumerate}

Estes termos foram buscados no título, abstract, e em palavras-chave nas bases que permitiam tal busca. Por exemplo, a string utilizada na IEEE Xplore foi:
\begin{verbatim}
("All Metadata":"gaming" OR "All Metadata":"game" OR 
 "All Metadata":"games") AND 
("All Metadata":"educational" OR "All Metadata":"serious") AND 
("All Metadata":"elderly" OR "All Metadata":"old adult" OR 
 "All Metadata":"old adults" OR "All Metadata":"senior")
\end{verbatim}

\subsection{Processo de Seleção}\label{subsec:processo_selecao}

O processo de seleção dos artigos foi realizado em quatro fases:
\begin{enumerate}
    \item \textbf{Fase 1}: Todos os artigos retornados pela pesquisa
    \item \textbf{Fase 2}: Remoção de artigos duplicados, escritos em outras línguas que não Inglês ou Português, ou outros tipos de trabalhos (definido pelo critério de exclusão 1)
    \item \textbf{Fase 3}: Leitura do título e \textit{abstract} e aplicação dos critérios de inclusão e exclusão
    \item \textbf{Fase 4}: Leitura do texto completo dos artigos, aplicação dos critérios de inclusão e exclusão (feita), e extração e análise dos dados relevantes (a ser realizada nas próximas etapas do PGC).
\end{enumerate}

Os detalhes do número de artigos encontrados e selecionados em cada fase são apresentados no Capítulo~\ref{cap:resultados}.

\section{Cronograma de Trabalho}\label{sec:cronograma}

O desenvolvimento deste projeto está organizado em três fases, correspondentes aos três quadrimestres do Projeto de Graduação em Computação (PGC-I, PGC-II e PGC-III), conforme detalhado na Tabela~\ref{tab:cronograma}.

\begin{table}[H]
\centering
\caption{Cronograma do Projeto}
\label{tab:cronograma}
\begin{tabular}{l|p{10cm}}
\hline
\textbf{Fase} & \textbf{Atividades Principais} \\ \hline
PGC-I (Fev–Abr/2025) & 
\begin{itemize}[leftmargin=*]
    \item Definição do tema e escopo da pesquisa
    \item Identificação de lacunas na literatura
    \item Execução das fases da RSL (busca, leitura e seleção dos artigos)
    \item Elaboração do relatório preliminar
\end{itemize} \\
PGC-II (Jun–Ago/2025) & 
\begin{itemize}[leftmargin=*]
    \item Leitura completa e análise dos artigos selecionados
    \item Extração e categorização das recomendações identificadas
    \item Elaboração de tabela comparativa das recomendações
    \item Elaboração do relatório intermediário
\end{itemize} \\
PGC-III (Set–Dez/2025) & 
\begin{itemize}[leftmargin=*]
    \item Consolidação dos resultados e análise
    \item Elaboração do relatório final
    \item Defesa do trabalho
\end{itemize} \\ \hline
\end{tabular}
\end{table}

\chapter{Resultados e Discussões}
\label{cap:resultados}

Ao término do processo de seleção, foram identificados 62 artigos elegíveis para a etapa de extração de recomendações e análise, conforme demonstrado na Tabela \ref{tab:selecao_artigos}. Esses estudos representam um panorama relevante de critérios elencados ou utilizados como referência no desenvolvimento de jogos digitais voltados ao público idoso, incluindo menções pontuais a aspectos relacionados à cibersegurança.

\begin{table}[H]
\centering
\caption{Número de artigos selecionados por fase da RSL}
\label{tab:selecao_artigos}
\begin{tabular}{lcccc|c}
\hline
\textbf{Fase} & \textbf{ACM} & \textbf{Web of Science} & \textbf{IEEE Xplore} & \textbf{SBC Open Lib} & \textbf{Total} \\ \hline
Fase 1        & 14           & 275                     & 167                  & 2                    & 458            \\
Fase 2        & 14           & 251                     & 156                  & 1                    & 422            \\
Fase 3        & 4            & 71                      & 36                   & 1                    & 112            \\
Fase 4        & 1            & 38                      & 22                   & 1                    & 62               \\ \hline
\end{tabular}
\end{table}

A listagem de cada um dos artigos e os critérios analisados em cada fase pode ser acessada em:
\begin{itemize}
    \item Fase 1: \href{http://bit.ly/3GtHj5X}{http://bit.ly/3GtHj5X}
    \item Fase 2: \href{https://bit.ly/4iYA4k1}{https://bit.ly/4iYA4k1}
    \item Fase 3: \href{https://bit.ly/3RIhUrp}{https://bit.ly/3RIhUrp}
    \item Fase 4: \href{https://bit.ly/4m14C7p}{https://bit.ly/4m14C7p}
    \item Resultado: \href{https://bit.ly/3YVzJqN}{https://bit.ly/3YVzJqN}
\end{itemize}

\section{Questões de Pesquisa}\label{sec:qp}

A realização dessas etapas da Revisão Sistemática da Literatura (RSL) possibilitou o início da resposta às questões de pesquisa propostas. A análise detalhada e a subsequente extração de dados dos artigos selecionados, a serem conduzidas nas próximas fases do Projeto de Graduação em Computação (PGC), fornecerão subsídios quantitativos para embasar respostas mais precisas e fundamentadas.

\subsection{Quais as recomendações, requisitos, e critérios identificados no desenvolvimento de jogos para o público idoso?}\label{subsec:qp1}

A maior parte dos estudos analisados apresenta o uso de jogos sérios com foco em objetivos como treinamento cognitivo e prevenção de doenças neurodegenerativas associadas à demência \cite{yang2024serious, zuo2024development, caggianese2018towards}, e promoção do envelhecimento ativo \cite{nacimiento-garcia2024gamification}.

Dentre as recomendações mais frequentemente mencionadas destacam-se: o uso de fontes ampliadas para facilitar a leitura \cite{tziraki2017designing}, a simplificação das interfaces visuais \cite{valladares2017design}, e a possibilidade de personalização de parâmetros como a velocidade do jogo, quando aplicável.

\subsection{Como jogos desenvolvidos tendo o público idoso como alvo são validados e testados?}\label{subsec:qp2}

Com base no levantamento inicial, observou-se que a validação de jogos voltados ao público idoso ocorre majoritariamente por meio da aplicação direta do jogo seguida de coleta de \textit{feedback} dos usuários \cite{merilampi2017cognitive}, pela aplicação de avaliações antes e após a intervenção \cite{wong2022effectiveness}, ou ainda por meio de consultas a especialistas das áreas da saúde e design, como fisioterapeutas, cuidadores e designers \cite{busca2024serious}.

\subsection{Que abordagens são usadas para levantamento de requisitos nestes jogos e quais produzem melhores resultados?}\label{subsec:qp3}

Verificou-se que poucos trabalhos descrevem de forma explícita a metodologia empregada para o levantamento de requisitos. Essa lacuna é abordada de forma específica em \cite{machado2018heuristics, manser2021making}.

Em \cite{machado2018heuristics}, foi conduzida uma revisão da literatura para levantamento de requisitos e princípios de design. Em seguida, jogos sérios para idosos oriundos da revisão da literatura foram testados com o público alvo, e avaliados por meio de um questionário, a fim de entender a validade de tais requisitos. Posteriormente, com base em dados empíricos, o estudo buscou desenvolver uma teoria que descreva características desejáveis no desenvolvimento de tais jogos.

Em \cite{manser2021making}, os autores executaram as seguintes fases: (1) Revisão da literatura. (2) Modelagem do perfil de usuário, definindo aspectos demográficos, capacidades, características, hobbies, e motivação para jogar. (3) Levantamento de necessidades terapêuticas junto a profissionais. (4) Definição de escopo tecnológico, ou seja, quais aparelhos eletrônicos seriam utilizados para o jogo. (5) Estratégia de manutenção, com o objetivo de possibilitar o uso da solução desenvolvida fora do período de estudo. O desenvolvimento do jogo se deu em um \textit{loop} de implementação dos requisitos seguido de validação com usuários.

\section{Desenvolvimento de jogos relacionados à cibersegurança}\label{sec:ciberseguranca}

Durante a etapa de busca, apenas \textbf{3} artigos atenderam ao critério de inclusão 4, que estabelece a necessidade de associação entre as recomendações identificadas e aspectos de cibersegurança ou proteção digital. Os trabalhos que cumpriram esse critério foram \cite{machado2017learning, bernardino2021serious, van2020serious}.

Nestes estudos, destacam-se aspectos como o estímulo ao pensamento crítico sobre segurança na internet por meio de jogos adaptados ao público idoso \cite{machado2017learning}, o uso de mecânicas lúdicas e narrativas cotidianas para promover a alfabetização digital em cibersegurança \cite{bernardino2021serious}, e a aplicação de agentes virtuais com análise vocal para treinar resiliência verbal contra fraudes praticadas presencialmente \cite{van2020serious}. Essas abordagens refletem diferentes estratégias para integrar conteúdos de proteção digital ao design de jogos sérios, com ênfase em situações reais enfrentadas por usuários vulneráveis.

A baixa quantidade de estudos que tratam simultaneamente de recomendações para jogos digitais e de princípios de cibersegurança evidencia uma lacuna significativa na literatura atual. Este cenário aponta para um campo de pesquisa ainda em estágio inicial, com amplo potencial para investigações futuras voltadas à integração entre segurança digital e design de jogos, especialmente para públicos vulneráveis, como o idoso.

\chapter{Conclusões}
\label{cap:conclusoes}

\section{Contribuições}

Este trabalho apresenta contribuições relevantes para o campo do desenvolvimento de jogos sérios voltados à cibersegurança para pessoas idosas. A principal contribuição consiste na realização de uma RSL que identifica e organiza critérios, recomendações e abordagens metodológicas aplicadas no design de jogos voltados para esse público, com ênfase em segurança digital. O estudo oferece uma visão abrangente do estado da arte, destacando lacunas importantes na literatura, especialmente no que se refere à ausência de diretrizes consolidadas e práticas participativas de levantamento de diretrizes sensíveis às necessidades cognitivas, físicas e tecnológicas dos idosos.

Além disso, o trabalho propõe um conjunto de recomendações que pode ser utilizado como ponto de partida para desenvolvedores e pesquisadores interessados em criar soluções mais acessíveis, inclusivas e eficazes. A sistematização de boas práticas encontradas nos estudos analisados contribui não apenas para a melhoria da qualidade dos jogos desenvolvidos, mas também para a promoção da inclusão digital segura desse grupo social cada vez mais conectado, porém vulnerável a ameaças cibernéticas.

Ao delimitar claramente os critérios de avaliação e exclusão de estudos, este trabalho também contribui metodologicamente para futuras revisões e pesquisas similares, fortalecendo a base teórica e empírica sobre o tema.

\section{Limitações}

Durante a etapa de busca dos artigos, observou-se uma predominância de estudos envolvendo jogos digitais que utilizam sensores de movimento, com destaque para os chamados \textit{exergames} — jogos desenvolvidos com o objetivo de promover a prática de atividades físicas. Considerando que esta revisão sistemática tem como foco a identificação de recomendações gerais para o desenvolvimento de jogos sérios que não dependam de aparatos tecnológicos específicos, como sensores externos, optou-se pela exclusão desses estudos por meio do critério de exclusão 6.

Além disso, constatou-se que muitos trabalhos não tinham como público-alvo principal a população idosa, mas sim indivíduos acometidos por condições clínicas frequentemente associadas ao envelhecimento. Como esta revisão busca compreender as necessidades específicas de pessoas idosas no contexto do desenvolvimento de jogos digitais, foram desconsiderados os estudos centrados em patologias específicas que não destacavam de forma clara a população idosa.

\section{Trabalhos Futuros}

Este estudo apresentou etapas iniciais da execução de uma Revisão Sistemática da Literatura (RSL) sobre o tema proposto. Nas próximas fases do Projeto de Graduação em Computação, conforme descrito na Seção~\ref{sec:cronograma}, está prevista a continuidade da RSL, incluindo a análise e discussão de novos resultados.

Além disso, propõe-se como trabalho futuro o desenvolvimento de um protótipo de jogo sério que incorpore as recomendações identificadas nesta revisão. A validação do protótipo deverá ser conduzida com a participação do público-alvo, bem como de cuidadores, especialistas em design de jogos e profissionais da área de cibersegurança, a fim de avaliar sua usabilidade, eficácia pedagógica e adequação às necessidades específicas dos usuários.



% ----------------------------------------------------------
% ELEMENTOS PÓS-TEXTUAIS (Referências, Glossário, Apêndices)
% ----------------------------------------------------------
\postextual

% Referências bibliográficas
\bibliography{bibliografia}


% Glossário (Consulte o manual)
%\glossary

% Apêndices
%% ----------------------------------------------------------
% Apêndices
% ----------------------------------------------------------

% ---
% Inicia os apêndices
% ---
\begin{apendicesenv}

% Imprime uma página indicando o início dos apêndices
\partapendices

% ----------------------------------------------------------
\chapter{Primeiro Apêndice}
% ----------------------------------------------------------


\end{apendicesenv}


% Anexos
%% ----------------------------------------------------------
% Apêndices
% ----------------------------------------------------------

% ---
% Inicia os anexos
% ---
\begin{anexosenv}

% Imprime uma página indicando o início dos anexos
\partanexos

% ---
\chapter{Nome do Primeiro Anexo}
% ---
\lipsum[30] % Texto qualquer. REMOVER!!

% ---
\chapter{Nome de Outro Anexo}
% ---

\lipsum[32] % Texto qualquer. REMOVER!!

\end{anexosenv}

% Índice remissivo (Consultar manual)
%\phantompart
%\printindex

\end{document}
